\section{Factorization over fields with solvable Galois closure}
\label{sec:solvable-factorization}

For $k\ge2$, recall the fixed-degree product
\begin{equation}
 \Gamma_k(T)=
 \prod_{\substack{\zeta^k=1\\ \zeta\ \mathrm{primitive}}}F_\zeta(T).
 \label{eq:Gamma-fiber-product}
\end{equation}

\begin{theorem}[Exact solvable-base factorization]
\label{thm:solvable-factorization}
Let $k\ge2$, put $K_k=\Q(\zeta_k)$, and let $B/\Q$ be any finite
extension whose Galois closure over $\Q$ is solvable.  Set
\[
 D=B\cap K_k.
\]
Then $\Gamma_k$ has exactly $[D:\Q]$ irreducible factors over $B$,
each of degree $6[K_k:D]$.  Equivalently, the factors are indexed by
the $\Gal(\overline\Q/B)$-orbits on the primitive $k$-th roots, and
the factor attached to an orbit $\mathcal O$ is
\[
 \Gamma_{\mathcal O}(T)=\prod_{\zeta\in\mathcal O}F_\zeta(T).
\]
In particular,
\[
 \Gamma_k\text{ is irreducible over }B
 \quad\Longleftrightarrow\quad B\cap K_k=\Q.
\]
\end{theorem}

\begin{proof}
Because $K_k/\Q$ is Galois, restriction identifies the image of
$\Gal(\overline\Q/B)$ in $\Gal(K_k/\Q)$ with
$\Gal(K_k/D)$.  Equivalently,
\[
 \Gal(BK_k/B)\cong\Gal(K_k/D).
\]
The action of $\Gal(K_k/\Q)$ on the primitive $k$-th roots is simply
transitive.  Hence there are $[D:\Q]$ relevant orbits, each of
cardinality $[K_k:D]$.

Let $\widetilde B/\Q$ be the Galois closure of $B/\Q$.  The extension
$\widetilde B K_k/K_k$ is finite Galois with solvable Galois group, so
\[
 BK_k\subseteq\widetilde B K_k\subseteq K_k^{\mathrm{sol}}.
\]
Corollary~\ref{cor:solvable-stability} therefore shows that every
$F_\zeta$ remains irreducible over $BK_k$.

Fix an orbit $\mathcal O$, choose $\zeta\in\mathcal O$, and let
$\alpha$ be a root of $F_\zeta$.  Since $F_\zeta(0)=-16$, one has
$\alpha\ne0$, and $\zeta=g(\alpha)$ gives
\[
 B(\alpha)=BK_k(\alpha).
\]
Consequently
\[
 [B(\alpha):B]
 =[BK_k(\alpha):BK_k]\,[BK_k:B]
 =6[K_k:D].
\]
The polynomial $\Gamma_{\mathcal O}$ belongs to $B[T]$, is monic,
has degree $6|\mathcal O|=6[K_k:D]$, and vanishes at $\alpha$.  It is
therefore the minimal polynomial of $\alpha$ over $B$.
Different orbit products are coprime: a common root $\alpha$ would
make $g(\alpha)$ belong to two disjoint orbits.  This proves that the
displayed products give the complete factorization.
\end{proof}

\begin{corollary}[Complete factorization over the maximal solvable extension]
\label{cor:Qsol-factorization}
For every $k\ge2$, equation \eqref{eq:Gamma-fiber-product} is the
complete irreducible factorization of $\Gamma_k$ over
$\Q^{\mathrm{sol}}$.  It consists of $\varphi(k)$ distinct sextics,
and the splitting field of each sextic over $\Q^{\mathrm{sol}}$ has
Galois group $A_6$.  The same product is also the complete
factorization over $\Q^{\mathrm{ab}}$.
\end{corollary}

\begin{proof}
Every primitive $k$-th root belongs to $\Q^{\mathrm{sol}}$, and
Corollary~\ref{cor:Qsol-monodromy} makes each $F_\zeta$ irreducible
over that field with splitting group $A_6$.  Distinct primitive roots
give distinct sextics because their $T^2$ coefficients are
$\zeta-5$.  Equation \eqref{eq:Gamma-fiber-product} proves
completeness.  Irreducibility over the larger field
$\Q^{\mathrm{sol}}$ also implies irreducibility over
$\Q^{\mathrm{ab}}\subseteq\Q^{\mathrm{sol}}$.
\end{proof}

\section{Solvable-base factorization of the delta factors}

Lemma~\ref{lem:root-correspondence} gives a
$\Gal(\overline\Q/\Q)$-equivariant bijection
\[
 \beta:Z(\Gamma_k)\longrightarrow
 Z(\widetilde\Delta_{4k,4}),
 \qquad \beta(\alpha)=h(\alpha).
\]

\begin{theorem}[Exact solvable-base delta factorization]
\label{thm:delta-solvable-factorization}
Let $k\ge2$ and let $B/\Q$ be a finite extension whose Galois closure
is solvable.  With $K_k=\Q(\zeta_k)$ and $D=B\cap K_k$, the polynomial
$\widetilde\Delta_{4k,4}$ has exactly $[D:\Q]$ irreducible factors
over $B$, each of degree $6[K_k:D]$.  Over $\Q^{\mathrm{sol}}$ its
complete factorization consists of $\varphi(k)$ irreducible sextic
factors, one for each primitive $k$-th root.  The splitting group of
each such sextic over $\Q^{\mathrm{sol}}$ is $A_6$.
\end{theorem}

\begin{proof}
Restricting the equivariance of $\beta$ to
$\Gal(\overline\Q/B)$ identifies its orbits, including their
cardinalities.  Theorem~\ref{thm:solvable-factorization} gives the
assertion over $B$, and Corollary~\ref{cor:Qsol-factorization} gives
the assertion over $\Q^{\mathrm{sol}}$.  On each sextic orbit the
two permutation representations are isomorphic, so their kernels and
hence their splitting fields agree; the group is therefore $A_6$.
\end{proof}

The same assertions hold for the unscaled factor
$\Delta_{4k,4}(c)$, because $C=4c$ is an invertible linear change over
$\Q$.

\begin{remark}
The preceding theorem is an orbit description; it does not identify
the sextic in the parameter variable with $F_\zeta(T)$.
\end{remark}

\begin{corollary}[Coprime cyclotomic base change]
\label{cor:coprime-cyclotomic-base}
Let $k\ge2$ and $\ell\ge1$, and let $\zeta_\ell$ be a primitive
$\ell$-th root of unity.  If $(k,\ell)=1$, then both $\Gamma_k$ and
$\widetilde\Delta_{4k,4}$ remain irreducible over
$\Q(\zeta_\ell)$.
\end{corollary}

\begin{proof}
For coprime $k$ and $\ell$,
$K_k\cap\Q(\zeta_\ell)=\Q$.  Apply Theorems
\ref{thm:solvable-factorization} and
\ref{thm:delta-solvable-factorization}.
\end{proof}

\begin{remark}
Corollary~\ref{cor:coprime-cyclotomic-base} is the period-$4$ analogue
of the period-$3$ result in
\cite[Theorem~1.2]{KoizumiMurakamiSanoTakehira2025}.  The theorems
above are stronger: they allow every finite base field with solvable
Galois closure and describe every resulting factor.
\end{remark}

\section{The solvable part of a parameter field}

\begin{corollary}[Maximal solvable subfield of a parameter field]
\label{cor:maximal-solvable-parameter-field}
Let $C=4c$ be an exact-period-$4$ parabolic parameter with primitive
$k$-th root multiplier $\zeta$, where $k\ge2$.  Then
\[
 \Q(C)\cap\Q^{\mathrm{sol}}=\Q(\zeta).
\]
Equivalently, $\Q(\zeta)$ is the largest subfield of $\Q(C)$ whose
Galois closure over $\Q$ is solvable.  In particular,
$C\notin\Q^{\mathrm{sol}}$.  The same assertions hold with $c$ in
place of $C$.
\end{corollary}

\begin{proof}
Let $\alpha$ be the corresponding root of $F_\zeta$.  The injective
Galois-equivariant root correspondence gives
\[
 \Q(C)=\Q(\alpha)=\Q(\zeta,\alpha).
\]
Put $K=\Q(\zeta)$ and $E=K(\alpha)$.  Corollary
\ref{cor:solvable-stability} gives
\[
 E\cap K^{\mathrm{sol}}=K.
\]
By Lemma~\ref{lem:solvable-closure-comparison},
$\Q^{\mathrm{sol}}\subseteq K^{\mathrm{sol}}$, while
$K\subseteq E\cap\Q^{\mathrm{sol}}$.  Therefore
\[
 K\subseteq E\cap\Q^{\mathrm{sol}}
 \subseteq E\cap K^{\mathrm{sol}}=K.
\]

If $M\subseteq E$ has solvable Galois closure over $\Q$, then
$M\subseteq\Q^{\mathrm{sol}}$ and hence $M\subseteq K$.  Conversely,
$K/\Q$ is abelian.  Thus $K$ is the largest such subfield.  If
$C\in\Q^{\mathrm{sol}}$, then $E=\Q(C)\subseteq\Q^{\mathrm{sol}}$,
contradicting $[E:K]=6$.  Finally $C=4c$, so the two coordinates
generate the same field.
\end{proof}

\begin{corollary}[Maximal Galois subfield]
\label{cor:maximal-galois-parameter-field}
Under the hypotheses of
Corollary~\ref{cor:maximal-solvable-parameter-field}, the cyclotomic
field $\Q(\zeta)$ is the largest subfield of $\Q(C)$ that is Galois
over $\Q$.
\end{corollary}

\begin{proof}
Write $K=\Q(\zeta)$ and $E=\Q(C)=K(\alpha)$.  The proof of
Theorem~\ref{thm:sextic-abelian-disjointness} gives the stronger fact
that $E$ contains no nontrivial intermediate field Galois over $K$.
If $D\subseteq E$ is Galois over $\Q$, then $DK$ is contained in $E$
and is Galois over $K$.  Hence $DK=K$, so $D\subseteq K$.  Conversely,
$K/\Q$ is Galois.
\end{proof}

At multiplier $1$, the exact-period-$4$ parameters instead satisfy
\[
 p(C)=C^3+9C^2+27C+135=0.
\]
After $U=C+3$, this is $U^3+108$.  The cubic is irreducible over
$\Q$, and its discriminant is
\[
 -27\cdot108^2=-2^4 3^9,
\]
which is not a square.  Its Galois closure has group $S_3$.
Consequently these parameters do not lie in $\Q^{\mathrm{ab}}$, but
they do lie in $\Q^{\mathrm{sol}}$.  This explains why the solvable
parameter-field statement above necessarily assumes $k\ge2$.

\section{Torsion Hilbert irreducibility context}

For this interpretation one must use the finite cover of degree six
\[
 g:U\longrightarrow\mathbf G_m,
 \qquad U=\mathbf G_m\setminus g^{-1}(0).
\]
Its degree is the degree of the rational function $g$, namely six.
After writing the target coordinate as $X=g(T)$, the element $T$
satisfies the monic equation $P_X(T)=0$ with nonzero constant term;
both $T$ and $T^{-1}$ are integral over
$\overline{\mathbf Q}[X,X^{-1}]$, which also proves finiteness.  The
pullback condition means that the fiber product with $x\mapsto x^n$
is geometrically irreducible for every $n\ge1$.  Its function-field
equation is
\[
 Y^n=g(T).
\]
Over $\overline{\mathbf Q}(T)$, the function $g$ is not a $p$-th
power for any prime $p$.  For $p\ge3$, this follows from its pole of
order two at $T=0$.  For $p=2$, suppose $g=w^2$.  Its only poles force
\[
 w=aT^2+bT+c+dT^{-1}.
\]
Comparison of the coefficients of $T^4,T^3,T^{-1},T^{-2}$ gives
\[
 a^2=-1,\qquad b=a,\qquad d^2=16,\qquad cd=4.
\]
Thus $c^2=1$, whereas comparison at $T^2$ gives
\[
 -1+2ac=-4,
 \qquad ac=-\frac32,
\]
so $(ac)^2=9/4$, contradicting $a^2c^2=-1$.  Hence $g$ is not a
square.

The field $\overline{\mathbf Q}(T)$ contains all roots of unity.  The
Kummer binomial criterion
\cite[Chapter~VI, \S9, Theorem~9.1]{Lang2002} says in this setting that
$Y^n-q$ is irreducible if and only if $q$ is not a $p$-th power for
every prime $p\mid n$.  In the general criterion one must also exclude
$q\in-4\overline{\mathbf Q}(T)^4$ when $4\mid n$; here every element
$-4b^4=(2ib^2)^2$ is already a square.  The preceding arguments
therefore prove that $Y^n-g(T)$ is irreducible for every $n\ge1$.

Zannier's cyclotomic Hilbert irreducibility theorem
\cite[Theorem~2.1]{Zannier2010} then gives finitely many exceptional
torsion fibers over the cyclotomic closure.  By Kronecker--Weber this
closure is $\Q^{\mathrm{ab}}$.  Corollary
\ref{cor:solvable-stability} determines the exceptions exactly: for
$\zeta\ne1$, the polynomial $F_\zeta$ is irreducible even over
$\Q(\zeta)^{\mathrm{sol}}$, while
\[
 F_1(T)=(T+2)(T^2+4)(T^3-2).
\]
Thus $\zeta=1$ is the unique exceptional torsion fiber for this cover.
